\scanpage{12}
(3.20)
\diagram{\resizebox{\linewidth}{!}{\begin{tikzcd}[ampersand replacement=\&,column sep=small,row sep=huge]
 X\arrow[d,"{\beta'}"'] \& = \& (X'\arrow[d,"{\textcircled{\scriptsize id}}"'] \& ,X''\arrow[d,"{\textcircled{\scriptsize $u''$}}"] \& ,X'\arrow[r,"u'"]\arrow[d,no head,"\mathrm{id}"] \& \varphi'(X'')\arrow[d,no head,"{\varphi'(u'')}"] \& ,X''\arrow[r,"{\deleted{\unsure{\varphi}}u''}"]\arrow[d,"u''"] \& \varphi''(X'))\arrow[d,no head,"\mathrm{id}"]\\
 g'f'(X) \& = \& (X' \& ,\varphi''(X') \& ,X'\arrow[r,"{\lambda'(X')}"] \& \varphi'\varphi''(X') \& ,\varphi''(X')\arrow[r,"\mathrm{id}"] \& \varphi''(X'))
\end{tikzcd}}}
\margin{adjonction}
et $\beta''$ analogue

[\uline{Question} Quand $\beta'$ est-il un \uline{épimorphisme} pour tout $X$ ? Cela \deleted{\unsure{implique|signifie}} \inserted{signifie}, moyennant des conditions d'exactitude : \unsure{droit|disons|droite} \unsure{standart|standard} sur $\varphi',\varphi''$, que $u'':X''\longrightarrow\varphi''(X')$ est tjrs un épimorphisme. ]

L'intersection des images essentielles de $g'$ et de $g''$ est formée des $(X',X'',u',u'')$ pour lesquels \deleted{\unsure{$g'$}}$u'$ \uline{et} \deleted{\unsure{$g''$}}$u''$ sont des isom. \deleted{Cette catégorie \unsure{i.}} Cette catégorie équivaut à l'image inverse par $g'$ de $g''_{\mathrm{ess}}(B'')$, qui est la sous-catégorie \inserted{str} pleine de $B'$ formée des $X'$ tels que $\lambda'(X')$ soit un isom. On trouve donc

(3.21) Conditions équivalentes :

(i) $g'_{\mathrm{ess}}(B')\wedge g''_{\mathrm{ess}}(B'')\simeq\text{catégorie finale}$

(ii) La sous-catégorie pleine de $B'$ formée des $X'$ tels que $\lambda'(X')$ soit un isom, est équivalente à la \deleted{sous}-catégorie \unsure{finale|final}

(ii bis) Idem pour $\lambda''$
% archive sheet number: 11
